\section{Coverage Threshold $\tau$ Robustness}
\label{app:tau-sweep}

The coverage thresholding recipe of \S\ref{sec:sib-design} is governed by a single global hyperparameter $\tau$ that sets the source-val error reduction criterion. The main-text reports use $\tau = 0.5$ (halve the source-val full-coverage error). This appendix verifies that the result is stable under perturbation of $\tau$.

\paragraph{Source-val operating-point verification.} Before testing on the held-out city, we confirm that the chosen $c^*_{val}$ achieves the $\tau$-criterion on the source-val split for every cell. Table~\ref{tab:tau-cstar-val} reports per-cell $c^*_{val}$, the source-val full-coverage error, the source-val error at $c^*_{val}$, and the achieved reduction percentage. All $9$ cells achieve at least the $\tau = 0.5$ target ($\geq 50\%$ reduction); reductions cluster slightly above $50\%$ ($50.1\%$ to $60.6\%$) because $c^*_{val}$ lands on the smallest grid value crossing the criterion. No cell falls back to the trivial $c = 1.0$.

\begin{table}[h]
\centering
\small
\setlength{\tabcolsep}{8pt}
\caption{\textbf{Source-val operating-point verification at $\tau = 0.5$.} For each cell: fitted $c^*_{val}$, source-val full-coverage error, source-val error at $c^*_{val}$, and achieved reduction.}
\label{tab:tau-cstar-val}
\begin{tabular}{llccccc}
\toprule
\textbf{Arch} & \textbf{Holdout} & \textit{$c^*_{val}$} & \textit{val err.\ at $c=1$} & \textit{val err.\ at $c^*_{val}$} & \textit{reduction} \\
\midrule
MAMNet  & Phoenix & $0.81$ & $0.239$ & $0.094$ & $60.6\%$ \\
MAMNet  & Miami   & $0.67$ & $0.362$ & $0.165$ & $54.5\%$ \\
MAMNet  & Chicago & $0.72$ & $0.334$ & $0.167$ & $50.1\%$ \\
OGLANet & Phoenix & $0.76$ & $0.315$ & $0.134$ & $57.6\%$ \\
OGLANet & Miami   & $0.62$ & $0.453$ & $0.204$ & $55.0\%$ \\
OGLANet & Chicago & $0.67$ & $0.399$ & $0.178$ & $55.5\%$ \\
DINOv3  & Phoenix & $0.72$ & $0.373$ & $0.163$ & $56.4\%$ \\
DINOv3  & Miami   & $0.57$ & $0.518$ & $0.257$ & $50.5\%$ \\
DINOv3  & Chicago & $0.62$ & $0.456$ & $0.221$ & $51.5\%$ \\
\bottomrule
\end{tabular}
\end{table}

\paragraph{$\tau$ sweep on the held-out city.} We re-fit $c^*_{val}$ at $\tau \in \{0.3, 0.5, 0.7\}$ on the source-val split and apply each fitted operating point to the held-out city. Table~\ref{tab:tau-sweep} reports the population-level test at each $\tau$. All three settings produce $9/9$ cells improving with Wilcoxon $p = 0.00195$ (the $n = 9$ floor under one-sided rank-sum), with the mean held-out reduction $\bar\Delta$ scaling monotonically with the criterion's aggressiveness. No cell falls back to the trivial coverage at any $\tau$.

\begin{table}[h]
\centering
\small
\setlength{\tabcolsep}{8pt}
\caption{\textbf{Population-level coverage recovery across $\tau$.} For each $\tau$: mean held-out selective shadow error reduction across $9$ cells, one-sided Wilcoxon $p$-value (alternative: $\Delta < 0$), cells improving, and median $c^*_{val}$.}
\label{tab:tau-sweep}
\begin{tabular}{ccccc}
\toprule
$\tau$ & \textit{Mean $\bar\Delta$ err.} & \textit{Wilcoxon $p$} & \textit{Cells improving} & \textit{Median $c^*_{val}$} \\
\midrule
$0.3$ & $-0.121$ & $0.00195$ & $9 / 9$ & $0.81$ \\
$0.5$ & $-0.199$ & $0.00195$ & $9 / 9$ & $0.67$ \\
$0.7$ & $-0.264$ & $0.00195$ & $9 / 9$ & $0.57$ \\
\bottomrule
\end{tabular}
\end{table}

\paragraph{Per-cell behavior across $\tau$.} Table~\ref{tab:tau-percell} reports per-cell $c^*_{val}$ and held-out selective shadow error reduction for each $\tau$. The pattern is consistent across cells: smaller $\tau$ yields a less aggressive operating point (larger $c^*_{val}$, smaller magnitude held-out reduction); larger $\tau$ yields a more aggressive operating point (smaller $c^*_{val}$, larger magnitude reduction). The DINOv3-Phoenix cell is the smallest in magnitude across all three settings, consistent with this cell's small within-city to LOCO gap (per Appendix~\ref{app:sib-comparison}); even on this cell, all $\tau$ settings produce a non-trivial reduction.

\begin{table}[h]
\centering
\scriptsize
\setlength{\tabcolsep}{4pt}
\caption{\textbf{Per-cell coverage recovery across $\tau$.} Held-out selective shadow error at $c^*_{val}$ minus error at $c = 1$, with $c^*_{val}$ in parentheses.}
\label{tab:tau-percell}
\begin{tabular}{ll|ccc}
\toprule
\textbf{Arch} & \textbf{Holdout} & $\tau = 0.3$ & $\tau = 0.5$ & $\tau = 0.7$ \\
\midrule
MAMNet  & Phoenix & $-0.046\ (0.91)$ & $-0.099\ (0.81)$ & $-0.128\ (0.76)$ \\
MAMNet  & Miami   & $-0.148\ (0.81)$ & $-0.230\ (0.67)$ & $-0.264\ (0.38)$ \\
MAMNet  & Chicago & $-0.115\ (0.86)$ & $-0.212\ (0.72)$ & $-0.271\ (0.24)$ \\
OGLANet & Phoenix & $-0.085\ (0.86)$ & $-0.131\ (0.76)$ & $-0.171\ (0.67)$ \\
OGLANet & Miami   & $-0.119\ (0.76)$ & $-0.228\ (0.62)$ & $-0.370\ (0.48)$ \\
OGLANet & Chicago & $-0.111\ (0.81)$ & $-0.226\ (0.67)$ & $-0.361\ (0.53)$ \\
DINOv3  & Phoenix & $-0.050\ (0.81)$ & $-0.085\ (0.72)$ & $-0.156\ (0.57)$ \\
DINOv3  & Miami   & $-0.225\ (0.72)$ & $-0.312\ (0.57)$ & $-0.349\ (0.34)$ \\
DINOv3  & Chicago & $-0.187\ (0.76)$ & $-0.268\ (0.62)$ & $-0.301\ (0.29)$ \\
\bottomrule
\end{tabular}
\end{table}

\paragraph{mIoU at threshold $0.5$ is preserved by construction.} Coverage thresholding does not modify the $0/1$ decision boundary at $p = 0.5$; it abstains on low-confidence pixels. The mIoU at threshold $0.5$ reported in Appendix~\ref{app:sib-comparison} is therefore identical for $\sib{}$ at $c = 1$ and at any $c^*_{val}$ in the tested range. The $\tau$ sweep characterizes the magnitude of selective-error recovery; deployment selects $\tau = 0.5$ as a single fixed global hyperparameter.